\documentclass[11pt,letterpaper]{article}
\usepackage[letterpaper,top=0.3in,bottom=0.3in,left=0.4in,right=0.4in]{geometry}
\usepackage[utf8]{inputenc}
\usepackage{mdwlist}
\usepackage[default]{lato}
\usepackage[T1]{fontenc}
\usepackage{textcomp}
\usepackage{fontawesome}
\usepackage{enumitem}
\usepackage{url}
\usepackage{hyperref}
\usepackage{needspace}

\hypersetup{hidelinks}

\DeclareFontFamily{U}{fontawesomeOne}{}
\DeclareFontShape{U}{fontawesomeOne}{m}{n}
{<-> FontAwesome--fontawesomeone}{}
\DeclareRobustCommand\FAone{\fontencoding{U}\fontfamily{fontawesomeOne}\selectfont}

\pagestyle{empty}
\setlength{\tabcolsep}{0em}

\newcommand{\headerrow}[2]
{\begin{tabular*}{\linewidth}{l@{\extracolsep{\fill}}r}
#1 & #2 \\
\end{tabular*}}

\begin{document}

\begin{center}
    {\LARGE \textbf{Minghao Guo}}
    \vspace{0.1cm}
    \\
    \raisebox{-0.2\height}{\Large \faPhoneSquare}\ (919) 306-3350
    \hfill \raisebox{-0.2\height}{\Large \faEnvelopeSquare}\ mhkwok2016@gmail.com
    \hfill \raisebox{-0.2\height}{\Large \faLinkedinSquare}\ \href{https://www.linkedin.com/in/minghao-guo-181b4b168/}{LinkedIn}
    \hfill \raisebox{-0.2\height}{\Large \faGraduationCap}\ \href{https://scholar.google.com/citations?user=T4w7A8IAAAAJ\&hl=zh-CN}{Google Scholar}
    \hfill \raisebox{-0.2\height}{\Large \faGlobe}\ \href{https://minghokwok.github.io/}{Homepage}
\end{center}

\hrule
\vspace{-1em}
\subsection*{\Large Profile}

\begin{itemize}[leftmargin=1em]
    \item Computer Science Ph.D. specializing in agentic retrieval, persistent memory, personalization, and evaluation. Builds systems that retrieve and reason over enterprise knowledge and long-term multimodal interactions, and evaluation frameworks that measure how personalized agents behave.
\end{itemize}

\hrule
\vspace{-1em}
\subsection*{\Large Education}

\begin{itemize}[leftmargin=1em]
    \parskip=0.1em

    \item
    \headerrow
        {\textbf{Rutgers University}}
        {\textbf{New Brunswick, NJ}}
    \headerrow
        {\emph{Ph.D. in Computer Science; Advisor: Dimitris N. Metaxas}}
        {\emph{Sep. 2023 -- Dec. 2026}}

    \item
    \headerrow
        {\textbf{North Carolina State University}}
        {\textbf{Raleigh, NC}}
    \headerrow
        {\emph{M.S. in Computer Science}}
        {\emph{Aug. 2021 -- May 2023}}

    \item
    \headerrow
        {\textbf{Sichuan University}}
        {\textbf{Chengdu, China}}
    \headerrow
        {\emph{B.E. in Computer Science and Technology}}
        {\emph{Sep. 2016 -- Jun. 2020}}
\end{itemize}

\hrule
\vspace{-1em}
\subsection*{\Large Experience}

\renewcommand\labelitemi{}
\renewcommand\labelitemii{$\bullet$}
\begin{itemize}[leftmargin=1em]
    \parskip=0.1em

    \item
    \headerrow
        {\textbf{Zoom}}
        {\textbf{Bellevue, WA}}
    \headerrow
        {\emph{Agentic Retrieval R\&D Intern}}
        {\emph{May 2026 -- Aug. 2026}}

    \Needspace{8\baselineskip}
    \item \emph{\textbf{Evaluation Framework for Context Memory of Zoom Meeting}}
    \begin{itemize}[leftmargin=1.3em,itemsep=0.05em,topsep=0.15em]
        \item Built an end-to-end evaluation framework for Zoom's meeting knowledge graph, generating synthetic meeting transcripts from event timelines with aligned ground-truth claims and QA pairs to measure retrieval quality and answer grounding.
        \item Designed a gated evaluation-and-improvement loop integrating extraction, clustering, retrieval, and claim-level scoring; improved claim recall from 0.625 to 0.750 on a representative evaluation set.
        \item Created a 260-question single-meeting benchmark covering action items, ownership, relationships, blockers, and temporal changes, enabling fine-grained diagnosis of retrieval and knowledge-capture failures.
    \end{itemize}

    \item
    \headerrow
        {\textbf{Rutgers University}}
        {\textbf{New Brunswick, NJ}}
    \headerrow
        {\emph{Ph.D. Researcher}}
        {\emph{Sep. 2023 -- Present}}

    \Needspace{6\baselineskip}
    \item \emph{\textbf{MemEye: A Visual-Centric Evaluation Framework for Multimodal Agent Memory}}
    \begin{itemize}[leftmargin=1.3em,itemsep=0.05em,topsep=0.15em]
        \item Developed a visual-centric benchmark for evaluating how multimodal agent memory systems retain, retrieve, and reason over information from long-term visual interactions.
        \item Proposed a two-dimensional evaluation taxonomy covering memory granularity, from scene-level context to pixel-level details, and evidence usage, from single observations to temporal state evolution.
        \item Evaluated 13 memory methods across four vision-language model backbones, revealing substantial weaknesses in fine-grained visual recall, temporal change tracking, and multi-observation evidence synthesis.
    \end{itemize}

    \Needspace{6\baselineskip}
    \item \emph{\textbf{Seeing, Maintaining, Learning: Evolution of Multimodal Agent Memory}}
    \begin{itemize}[leftmargin=1.3em,itemsep=0.05em,topsep=0.15em]
        \item Led 80 authors and coordinated taxonomy design and research synthesis, organizing methods by experience representation and the evolution of memory-management policies, from static retrieval to self evolving.
        \item Analyzed 252 memory architectures and 76 benchmark records, identifying gaps in evidence preservation, temporal consistency, and bounded lifelong memory.
    \end{itemize}

    \Needspace{6\baselineskip}
    \item \emph{\textbf{DeepSieve: Information Sieving via LLM-as-a-Knowledge-Router}}
    \begin{itemize}[leftmargin=1.3em,itemsep=0.05em,topsep=0.15em]
        \item Built an agentic RAG framework that decomposes complex questions, routes subquestions to documents, structured data, and APIs, and iteratively filters retrieved evidence.
        \item Enabled reasoning across heterogeneous sources without merging their indexes; evaluated multi-hop QA on MuSiQue, 2WikiMultiHopQA, and HotpotQA.
    \end{itemize}

    \Needspace{6\baselineskip}
    \item \emph{\textbf{Individual Turing Test: Personalized Agent via RAG, LoRA \& Hybrid Memory}}
    \begin{itemize}[leftmargin=1.3em,itemsep=0.05em,topsep=0.15em]
        \item Built a personalized agent by fine-tuning Qwen2.5-7B with LoRA on longitudinal personal WeChat chat data, and compared fine-tuning, RAG, and hybrid memory through blind judgments by the user's acquaintances.
        \item LoRA fine-tuning alone beat a RAG-only baseline in blind pairwise judgments (win rate 0.608, p $<$ 0.001), and combining RAG with LoRA performed best, showing personalization benefits from both memory retrieval and model adaptation.
        \item Proposed an individual Turing test framework where acquaintances evaluate the authenticity of the AI digital twin against the real user.
    \end{itemize}

    \Needspace{8\baselineskip}
    \item
    \headerrow
        {\textbf{Tsinghua University}}
        {\textbf{Beijing, China}}
    \headerrow
        {\emph{Ph.D. Research Intern}}
        {\emph{May 2025 -- Aug. 2025}}

    \item \emph{\textbf{TwinVoice: A Multi-dimensional Benchmark Towards Digital Twins via LLM Persona Simulation}}
    \begin{itemize}[leftmargin=1.3em,itemsep=0.05em,topsep=0.15em]
        \item Designing \textbf{TwinVoice}, the first \textbf{real-world} benchmark to evaluate Digital Twin agents (LLM-based systems that emulate a target human’s linguistic style, preferences, and personality).
        \item Curating large-scale datasets across \textbf{three tracks}: Social Behavior, Personal Dialogue, and Narrative Persona, covering both \textbf{discriminative} and \textbf{generative} tasks.
        \item Leveraging \textbf{authentic conversational history} to assess style imitation fidelity, memory consistency, and persona alignment, avoiding synthetic or crowd-annotated biases.
    \end{itemize}

\end{itemize}

\hrule
\vspace{-1em}
\subsection*{\Large Selected Publications}

\begin{enumerate}[leftmargin=1.6em,itemsep=0.15em,parsep=0pt,topsep=0.1em]
    \item \textbf{MemEye: A Visual-Centric Evaluation Framework for Multimodal Agent Memory.} \\
    \textbf{Minghao Guo}*, Qingyue Jiao*, et al. \\
    \textit{Under review at \textbf{ICLR 2027}.}\href{https://arxiv.org/abs/2605.15128}{ [Paper]}

    \item \textbf{Seeing, Maintaining, and Learning: The Evolution of Multimodal Agent Memory.} \\
    \textbf{Minghao Guo}* (project organizer), Qingyue Jiao*, et al. \\
    \textit{Under review at \textbf{TMLR}.}\href{https://arxiv.org/pdf/2610.19366}{ [Paper]}\href{https://github.com/Seeing-Maintaining-Learning}{ [Project \& Resources]}

    \item \textbf{DeepSieve: Information Sieving via LLM-as-a-Knowledge-Router.} \\
    \textbf{Minghao Guo}, Qingcheng Zeng, et al. \\
    \textit{\textbf{Findings of EACL 2026}.}\href{https://arxiv.org/abs/2507.22050}{ [Paper]}\href{https://github.com/MinghoKwok/DeepSieve}{ [Code]}

    \item \textbf{Individual Turing Test: A Case Study of LLM-based Simulation Using Longitudinal Personal Data.} \\
    \textbf{Minghao Guo}, Ziyi Ye, et al. \\
    \textit{\textbf{SIGIR 2026}, short paper.}\href{https://arxiv.org/abs/2603.01289}{ [Paper]}

    \item \textbf{TwinVoice: A Multi-dimensional Benchmark Towards Digital Twins via LLM Persona Simulation.} \\
    Bangde Du*, \textbf{Minghao Guo}*, et al. \\
    \textit{\textbf{ACL 2026}.}\href{https://arxiv.org/pdf/2510.25536}{ [Paper]}
    \item \textbf{OpenPM: Auditable Point-in-Time Evaluation for LLM Portfolio-Management Agents.} \\
    Xinying Cai*, \textbf{Minghao Guo}*, et al. \\
    \textit{Forthcoming in \textbf{SIGKDD Explorations}, Dec. 2026.}\href{https://arxiv.org/abs/2608.09988}{ [Paper]}


    \item \textbf{ReaGAN: Node-as-Agent-Reasoning Graph Agentic Network.} \\
    \textbf{Minghao Guo}, Xi Zhu, et al. \\
    \textit{\textbf{HCRS @ TheWebConf 2026}.}\href{https://arxiv.org/pdf/2508.00429}{ [Paper]}
\end{enumerate}

\noindent\textit{* Equal contribution.}

\hrule
\vspace{-1em}
\subsection*{\Large Skills}

\begin{itemize}[leftmargin=0.15in]
    \item \textbf{Programming \& Systems:} Python, SQL, C++, Bash, Git, Docker, Linux
    \item \textbf{ML \& Retrieval:} PyTorch, Hugging Face, vLLM, RAG, vector search, knowledge graphs, LLM agents, post-training (SFT, LoRA)
    \item \textbf{Evaluation:} benchmark design, retrieval and claim-level evaluation, multi-hop QA, personalized-agent evaluation
\end{itemize}

\end{document}
